\chapter{予備審査会における指摘事項への回答と補足説明}
\label{chap:qa_discussion}

本章では，修士論文予備審査会において指摘された事項および質問に対し，本研究の立場から整理した回答と補足的な考察を示す．特に，LIPM における CoM 高さの位置づけ，軟弱地面における沈み込み補償の意味，本手法の有効範囲，ならびに既存研究との関係性について明確化することを目的とする．

\section{LIPM において CoM 高さを維持することの意味}
\label{sec:qa_com_height_meaning}

LIPM（Linear Inverted Pendulum Model）に基づく歩行制御では，重心（CoM）の高さを一定に保つことが基本的な仮定として用いられる．これは，CoM 高さを一定とすることで，ZMP と CoM の水平運動の関係が線形化され，制御設計が著しく簡略化されるためである．

LIPM における水平方向運動は
\begin{equation}
\ddot{x} = \frac{g}{z_{\mathrm{c}}}(x - x_{\mathrm{zmp}})
\end{equation}
で記述され，この式から分かるように，CoM 高さ $z_{\mathrm{c}}$ は運動の自然周波数を決定する重要なパラメータである．CoM 高さが一定であることで，ZMP 軌道設計と CoM 運動制御の対応関係が時間的に安定し，歩行生成および安定化制御が成立する．

したがって，LIPM において CoM 高さを「目標として維持する」という考え方は，物理的な意味だけでなく，モデルの前提条件を満たし続けるという制御上の意味を持つ．

\section{軟弱地面における CoM 高さ低下とその補償の位置づけ}
\label{sec:qa_sink_compensation}

軟弱地面を踏んだ場合，足裏が沈み込むことで，実環境における CoM 高さは計画時よりも低下する．このとき，制御系が参照する CoM 高さと実際の CoM 高さの間に乖離が生じ，LIPM の仮定が破綻する．

本研究における CoM 高さ制御は，「常に CoM を高くする」ことを目的とするものではなく，「沈み込みによって低下した CoM 高さを補償し，LIPM の前提条件に近い状態を回復させる」ことを目的としている．すなわち，本手法は沈み込みを検知した局面においてのみ CoM 高さ参照を一時的に引き上げ，DS--SS 遷移時の横方向倒れ込みおよび支持脚足首トルクの増大を抑制するものである．

この観点から，本研究における CoM 高さ制御は，
「軟弱地面への踏み込みによって実効的な CoM 高さが低下し，
LIPM が仮定する力学条件から逸脱する状況に対して，
その逸脱を補償する制御」であると位置づけられる．

\section{単純な変更であるにもかかわらず効果が大きい理由}
\label{sec:qa_simple_but_effective}

本研究で提案した手法は，既存の LIPM ベース制御器の構造を大きく変更するものではなく，CoM 高さ参照の更新則に対する比較的単純な修正に基づいている．しかしながら，シミュレーション結果から示されたように，DS--SS 遷移時の支持脚ロール方向トルクや ZMP の発散挙動に対して顕著な改善効果を示した．

この理由は，DS--SS 遷移という不安定化が集中する局面において，CoM 高さが倒立振子の角加速度を直接支配するパラメータであるためである．すなわち，制御全体を複雑化せずとも，不安定化の主因に直結する変数を適切に操作することで，大きな安定化効果が得られる．

この点は，本研究の重要な特徴であり，「シンプルであるがゆえに既存制御器との親和性が高い」という利点にもつながっている．

\section{提案手法の適用範囲と限界}
\label{sec:qa_applicability_limit}

一方で，本研究の結果から，提案手法がすべての軟弱地面条件に対して有効であるわけではないことも明らかとなった．特に Case3 に示されるように，沈み込み速度および沈み込み量が極めて大きい条件では，CoM 高さの調整以前に姿勢の不安定化が急激に進行し，転倒を回避することが困難であった．

このことから，本研究の提案手法には明確な適用範囲が存在し，「沈み込みによる不安定化が CoM 高さ制御によって緩和可能な範囲」において有効であると整理できる．常時軟弱な地面や極端な沈み込みを伴う環境に対しては，脚剛性制御や接触モデルの変更など，別のアプローチとの併用が必要になると考えられる．

\section{グラフの読み方に関する補足説明}
\label{sec:qa_graph_explanation}

本研究では，ZMP 軌道，CoM 高さ，支持脚ロール方向トルクの時間履歴を用いて歩行安定性を評価している．ZMP 軌道については，支持多角形内に留まっているか，あるいは外側へ発散しているかが安定性の判断基準となる．

CoM 高さのグラフは，沈み込み発生後で実環境の CoM がどの程度影響し，それに対して参照値がどのように補償されているかを示している．支持脚ロール方向トルクの時間履歴は，DS--SS 遷移時に足首関節にどれだけの力学的負荷が集中しているかを可視化したものであり，ピーク値の低減が安定化効果を示す指標となる．

これらのグラフを併せて解釈することで，「ZMP が発散しない理由」「トルクが低減された理由」「歩行が継続できた要因」を段階的に理解できる構成となっている．

\section{軟弱地面に対する既存アプローチとの比較}
\label{sec:qa_related_approaches}

軟弱地面への適応に関する既存研究としては，脚剛性や減衰をモデルパラメータとして扱い，それらを適応的に調整するモデルベース手法が数多く提案されている．例えば，Dual-SLIP モデルを用いた研究では，床の不確かさを含むダイナミクスを陽に扱い，ロバスト性を確保する制御設計が行われている．

これに対し，本研究は LIPM を基盤とし，既存の ZMP ベース歩行制御との整合性を保ったまま，CoM 高さという単一の変数を操作することで沈み込み床への適応を実現する点に特徴がある．すなわち，モデルの高度化によって問題を解決するのではなく，「どの局面で何が破綻しているか」を解析し，その破綻点に直接作用する制御を付加するという立場を取っている．

この点において，本研究は軟弱地面適応に対する別の実用的アプローチとして位置づけられる．
